\documentclass[conference]{IEEEtran}


\usepackage{cite}
\usepackage{pbox}
\usepackage{array}
\newcolumntype{P}[1]{>{\centering\arraybackslash}p{#1}}
\newcolumntype{M}[1]{>{\centering\arraybackslash}m{#1}}
\usepackage{amsmath}
\usepackage{algorithmic}
\usepackage{array}
\usepackage{graphicx}
% \usepackage{subfigure}
% \usepackage{subcaption}
\usepackage{tikz}
\usetikzlibrary{arrows, calc, shapes, fit, positioning}
% \usepackage{caption}
\usepackage{pgfplots}
\pgfplotsset{compat=1.12}
\usepackage{xr}
%\usepackage{mathtools}
\usepackage{amsmath}

\usepackage{url}
\usepackage[utf8]{inputenc} 
\usepackage{fourier} 

% \usepackage{enumitem}
\usepackage[inline, shortlabels]{enumitem}
% \hyphenation{op-tical net-works semi-conduc-tor}

% http://tex.stackexchange.com/questions/56529/note-below-figure
\usepackage[capposition=top]{floatrow}

% \usepackage{forest}
 \usepackage[edges]{forest}
 
\usetikzlibrary{shapes.geometric}
% \usepackage[T1]{fontenc}

\usepackage{caption}
\usepackage{subcaption}

\begin{document}
%
% paper title
% Titles are generally capitalized except for words such as a, an, and, as,
% at, but, by, for, in, nor, of, on, or, the, to and up, which are usually
% not capitalized unless they are the first or last word of the title.
% Linebreaks \\ can be used within to get better formatting as desired.
% Do not put math or special symbols in the title.

% \title{Enhance the Mobility of Edge Servers Based on Live Migration of Docker Containers}

% \title{Offloading Service Handoff Across Edge Servers Through Live Migration of Docker Containers}

\title{Offloading Service Handoff Across Edge Servers Through Live Migration of Docker Containers}

% author names and affiliations
% use a multiple column layout for up to three different
% affiliations
\author{\IEEEauthorblockN{ Lele Ma, Qun Li}
% \IEEEauthorblockN{ Qun Li}
\IEEEauthorblockA{The College of William and Mary\\
Email: \{lma03, liqun\}@email.wm.edu}
}



% make the title area
\maketitle


\begin{abstract}

Supporting free mobility of mobile clients is vital important for offloading service on edge computing platform. The free mobility demands seamless migration of the offloading service among its nearby edge servers.
% in order to keep such promise
% .
% 
However, fast migration of offloading services across edge servers under WAN network environment poses significant challenges to the design of service handoff. 
%
%
%
In this paper, we present a novel service handoff system, which seamlessly migrates offloading service to the client's nearest edge server while the mobile client is moving. 
% Offloading services are running inside Docker containers. 
Service handoff is achieved via container's migration. 
% 
We identify the performance problem of the state-of-the-art work for Docker container's migration.
% We analyze storage management of Docker containers based with AUFS storage driver .
%
Base on our systematically study of the container layer management and images stacking,
we propose a migration method which leverages the layered storage system to reduce file system synchronization overhead without depending on the distributed file system. 
We implement a prototype system and conduct experiments using real world product applications.
Evaluation results reveal that, compared to the state-of-the-art service handoff system designed for edge computing platform, our system reduces the total duration of service handoff time by 
$80\%$($56\%$) with network bandwidth as 5Mbps(20Mbps).



% Edge computing platform provides low end-to-end latency by offloading mobile device's resource-intensive processing to the edge servers.
% Edge server provides low end-to-end latency of resource-intensive application by providing offloading services for its nearby mobile clients.
% Offloading mobile device's workloads to its nearby edge server provides high responsive user experiences for the client. 

% Offloading service provided by edge servers should support the free mobility of the mobile clients. 
% While edge computing platform promises to provide low end-to-end latency services to its nearby mobile clients, 
% the support of client mobility has great impacts on the service performance. 



% Specifically, we aim to reduce the total duration of migration instead of the downtime.
%
%
%
% This paper provides a sample of a \LaTeX\ document which conforms,
% somewhat loosely, to the formatting guidelines for
% % ACM SIG Proceedings.\footnote{This is an abstract footnote}

% Edge computing platform provides low end-to-end latency by offloading mobile device's resource-intensive processing to the edge servers.
% Edge server provides low end-to-end latency of resource-intensive application by providing offloading services for its nearby mobile clients.
% Offloading mobile device's workloads to its nearby edge server provides high responsive user experiences for the client. 
\end{abstract}





\section{Introduction}

Edge computing has attracted lots of attention both from the industry and academia in recent years\cite{satya2017edge}, \cite{yi2015fog,yi2015survey},\cite{satya2009case}. By placing the computing and storage nodes in close proximity to mobile devices or sensors, it enables highly responsive and reliable cloud services for mobile computing\cite{satya2017edge}. 

Offloading the computation to the nearest edge servers could reduce network latency and improve user experiences. When a mobile device moves away from its current edge server, the network latency will increase, which will highly reduce the performance of offloading services. Ideally, when the user moves, its offloading service running on the edge server should also move to another edge server in order to get low network latency. Therefore, migrate the computation lively from one edge server to a nearer edge server is vital important for the edge computing infrastructure. 

VM Hand-off \cite{ha2015vmhandoff} has been proposed to accelerate the service hand-off across the offloading edge servers. It divided VM images into two stacked overlays based on Virtual Machine (VM) synthesis \cite{satya2009case} techniques. So that the mobile device only need to transfer the VM's top application overlay to the target server instead of the whole VM image volume. However, the total transferred size is still in the magnitude of tens of megabytes and even hundreds of megabytes. The total handoff time was $1\sim4$ minutes on a WAN network, which is still a relatively long time for latency sensitive mobile applications. Furthermore, VM image overlays is hard to maintain and not widely available due to its limited support and deployment in the real world.

However, the wide deployment of Docker platform enlightens the possibility of high speed offloading service hand-off. As a container engine, Docker is increasingly popular in industrial cloud platform. It serves as an composing engine for Linux containers, where applications are running in an isolated environment based on the OS-level virtualization technique.
Docker's application runs in an isolated container with all its dependencies installed. Docker's storage driver uses layered images inside containers, which enables fast packaging and shipping of any application as a lightweight container. 
% With Docker platform, the application was encapsulated into containers and the storage system has inherently layered images.  

Docker's layered storage inherently supports copy-on-write where application's all writable data are encapsulated into one thin layer on top of its base image layers. This is also an excellent support for the high speed migration of services in this paper. Based on the layered storage, we can only transfer the thin writable layer on top of the container's layer stacks during migration. The application's base image layers which is usually large in size can be downloaded before the migration. This is applicable because Docker infrastructure inherently uses the cloud storage of container images (like in DockerHub), where all the container images are available anywhere across the Internet. 

But Docker itself is still in active development, and it doesn't supports live migration of its containers yet. So in order to enable offloading service hand-off based on Docker platform, we must meet the following challenges.

Firstly, we need to know the details of how Docker manages its image layers and how a container is created from layerred images. Until now, we cannot see any research work that introduce the images management inside Docker. In order to provide a systematic bird-view over Docker's storage management, we investigate and summarize the layered images based on AUFS storage driver in \ref{aufsIntroduction}.

Secondly, we need to find out how can we make use of Docker's layered images to speed up the live migration of containers. According what we found inside AUFS storage, Docker will create a new SHA256 number as local  identification for each image layer downloaded from the cloud. 
This means that if two Docker host download the same image layer from the same repository, they will have different reference identification. This is originally a safety mechanism to avoid the image overlapping across Docker hosts\cite{dockerlayer}. However, when we migrate one container from one Docker host to another Docker host, we must recognize those image layers that have different local identifications but actually have the same content with the same origin. So that those image layers will not be transferred during the migration.

Thirdly, the transferred data will include both the file system as well as the checkpointed binary memory images. Although the total size is reduced dramatically by leveraging Docker's layered storage, we still need to find a way to further reduce the transfer size in order to reduce the transfer time. 

Last but not the least, we need to find a way to mitigate the user-experience latency caused by migrating a container across different edge servers. Use experienced latency could be shorter than the actual migration time by well designed strategies of migration process. Ideally, we should make it close to a seamless service hand-off where users cannot notice their offloading edge server has been changed.

In order to meet those challenges, we proposed a framework that could enable high speed offloading service hand-off across edge servers by only encapsulate and transfer the thin writable layer of the application and its runtime status based on Docker's  experimental feature of checkpoint and restore. The contributions of this paper are listed as below:

\begin{itemize}
    
    \item It analyses the status of current techniques to do live migration for Docker containers. It evaluates the performance of some state-of-the-art works that could do Docker container live migration and expose the problems in it.
    
     \item It analyses the detailed storage management of Docker containers based on AUFS storage driver. It systematically introduces the layer management and images stacking methodology inside the Docker storage and reveals the proper way to leverage those mechanisms to speed up the live migration.
    
    
    \item It proposes a framework that enable service hand-off across edge offloading servers with help of the high speed live migration of containers leveraging the layered images of Docker platform. The framework ensures low end-to-end latency of resource-intensive cloud offloading while maintaining the high mobility of clients. 
    
    \item It implemented a prototype of the system and evaluate the hand-off performance of containers and response time of offloading applications, like OpenFace, etc. 
    
\end{itemize}

\section{Background and Related Work}

In this section, we further discuss the motivations by answering
questions: why edge applications need computation offloading, why we need  
service handoff in edge computing, why we need migration for service handoff, and why we want to do service handoff via container migration.

\subsection{Emerging Applications in Edge Computing Call For Computation Offloading}

With the rapid development of edge computing, many researchers have built various applications to take the advantages of the edge computing platforms. 
One example of such services is Augmented Reality (AR), where the application on a mobile device overlays augmented reality content
onto objects viewed on the device's camera and edge servers can
provide local object tracking and local AR content caching
\cite{satya2009case,MEC2014initiative,MEC2015-5G,hao2017challenges}.
% The solution minimizes round trip time and maximizes throughput for optimum quality of experience. 
% It can be used to offer consumer or enterprise propositions, such as tourist information, sporting event information, advertisements etc.
%
% Sataya Gabriel 
Gabriel platform \cite{ha2014wearable}  was proposed within the contexts of wearable cognitive assistance applications using a Glass-like wearable device, such as Lego Assistant, Drawing Assistant, Ping-pong Assistant.%,  and Assistance from Crowd-sourced Videos. 
Those applications need to react to a user's actions in real time with predefined guidance or guidance from crowd-sourced videos. 
%
%
% openface
OpenFace\cite{openface2016} is a real-time face recognition application for mobile that offers high accuracy with low training and prediction times. Mobile clients send the captured pictures from the camera to the nearby server. The server runs face recognition task with deep neural networks and send symbolic feedback to mobile clients in real time.  
More edge applications can be found in \cite{yi2015fog,yi2015survey,satya2017edge}.
%
All those edge applications need to offload intensive computations (e.g. computer vision, signal processing) to the edge server in order to get real time response. However, there are challenges before we can pragmatically deploy computation offloading in edge computing architectures. 

\subsection{Effective Edge Offloading Needs Seamless Handoff}

As we have mentioned, most edge applications can greatly benefit from offloading its heavy operators to the edge. However, the real time responsive service largely relies on the relative near distance to the edge servers under a WAN network. When the mobile client moves far away increasing the its distance to the current edge server, offloading performance bonus will be diminished dramatically.
Therefore, effective edge offloading needs mobility support.

In the centralized cloud infrastructure, mobility of clients can be well supported since the cloud server is always connected by clients through the long latency WAN internet. However, in the edge computing infrastructure, mobile devices are connected to the nearby edge server within its limited network coverage scope. Therefore, when the mobile device moves out of the edge server's coverage, it might be out of service, or alternatively, it has to ``handoff'' from the current edge server to a nearer edge server. 
This is an analogy to the \textit{seamless handoff} or \textit{handover}  mechanism in the cellular networks, where the moving client can seamlessly handoff across base stations and be served by its nearest available station. 
%
Therefore one of the primary needs for edge computing is to enable the free mobility for the offloading computations on the edge servers along with the mobile devices via seamless handoff.


\subsection{Seamless Edge Handoff via VM Migration}

However, there exists one key difference between the cellular network handoff and edge server handoff.
%
In cellular networks, changing a base station for a client is as simple as rebuilding a wireless connection. Most of the run-time states of the services are still available either on the client or on the remote server. After re-connection, the run-time state can be quickly seamlessly resumed. 
%
In the edge infrastructure, most mobile devices are using edge server to offload the resource-hungry or computation-intensive computations. This means that the edge server needs to hold the states of all those resource intensive workloads. 
In our mobile handoff from one edge server to another, just rebuilding the connection is certainly not enough. Since all the offloaded workloads are also need to be transferred from the current edge server to the next edge server.  

One possible solution is to use the live migration\cite{vmlivemig} to migrate a virtual machine from on edge server to another in order to seamlessly transfer the offloading workloads. However, this approach is already shown to be not suitable for edge computing environment in  \cite{ha2015vmhandoff}. Firstly, live migration and service handoff are optimized for very different performance metrics. While live migration aims to reduce the \textit{downtime} of the virtual machine, service handoff aims to reduce the \textit{total time} to the completion upon the handoff request is received. This is well discussed in \cite{ha2015vmhandoff}.
Secondly, live migration is originally designed for high performance data centers with high bandwidth network. However, this is not possible for edge servers which are deployed on the edge network over the WAN. Furthermore, live migration relies on shared storage so that it transfers only run-time state and doesn't transfer the disk storage data. Apparently, shared storage across the edge computing infrastructure is not feasible due to its widely distributed nature and low bandwidth across WAN. 

In order to enable handoff across edge computing servers, many more researches have been done based on virtual machine ~\cite{satya2009case,ha2015vmhandoff}. 
However, the total handoff time was still several minutes on a WAN network. For example, it will requires $245.5$ seconds to migrate a running OpenFace instance under 5Mbps bandwidth (50ms latency) network in \cite{ha2015vmhandoff}. 

One of the reasons for the long latency of handoff the large image size we need to transfer. Although the VM synthesis could reduce the image size by splitting images into multiply layers and only transfer the application-specific layer, the total transferred size is still in the magnitude of tens of megabytes and even hundreds of megabytes. 
The application layer was encapsulated with the whole application, both its static binary programs and its runtime memory data, which we think is an unnecessary cost. 

On the other hand, the deployment of the VM synthesis system is challenging for the legacy system. In order to enable VM synthesis, the hypervisor of VMs need to be patched to track the dirty memory at runtime. Also the storage of VM images need to be adapted to Linux FUSE interfaces in order to track file system changes inside the running VMs. Those two changes are hard to deploy in practice since they changes the behavior of hypervisors and file systems of the legacy platform and will also cause lots of performance overhead. 


\subsection{Why We Need Migration of Docker Container}
% Since handoff approaches base on VM migration posts significant performance pressure for seamless handoff of edge services, the container-based live migration has gained attractions in building efficient edge computing systems. % here we can cite quite a lot papers.
% As a container engine, Docker is increasingly popular in industrial cloud platform. 
Docker\footnote{Through out paper, we use Docker as our container engine implicitly.} is an composing engine for Linux containers, where applications are running in an isolated environment based on OS-level virtualization. Docker enables layered storage inside containers, which allows fast packaging and shipping of any application as a lightweight container. Each Docker image references a list of read-only layers that represent filesystem differences. Layers are stacked on top of each other to form a base for a container's root filesystem \cite{dockerlayer}. 

This layered storage allows fast migration of containers' run-time states without having to transfer the system and application base images. With the cloud storage of container images (like in DockerHub), all the container images are available anywhere across the Internet. Therefore, before migration starts, any edge server has the chance to download the system and application images as the base images stack for the container to run on. 
During the migration, we only need to transfer the run-time memory states and the thin \textit{container layer} on top of the Docker images stack. 

Apparently, the migration of the Docker containers is a better choice than the virtual machine based approaches we introduced above. The layered storage in Docker infrastructure enlightens a great opportunity for the service handoff based on the container migration. If the migration of containers can be done very fast, we can have the near seamless offloading service handoff across the adjacent edge servers.

However, at the time this paper is written, there is no tool that can leverage Docker's layered images to migrate containers efficiently. 
% In order to do this, we must fist dive into the images layers underlying containers and avoid the transfer of application images stack. That is just transfer the ``\textit{container layer}''. 

\section{Design and Implementation}



% % \subfigcapmargin = .5cm
% \begin{figure}
% \begin{subfigure}[b]{0.4\textwidth}
%     % \centering
%     \includegraphics[width=\linewidth]
%     % {figure/test-tarssh.eps}
%     {figure/sys.png}
%     \caption{Before Migration}
%     \label{subfig:before}
% \end{subfigure}
% \quad\quad\quad\quad
% \begin{subfigure}[b]{0.4\textwidth}
%     % \centering
%     \includegraphics[width=\linewidth]
%     % {figure/test-tarssh.eps}
%     {figure/sys2.png}
%     \caption{After Migration}
%     \label{subfig:after}
% \end{subfigure}
% \caption{ Offloading Serivce Handoff: Before and After Live Migration of Offloading Container. }
% \label{fig:before-after}
% \end{figure}

% https://tex.stackexchange.com/questions/350729/software-architecture-picture-in-tikz
\tikzset{
    srv/.style  = {draw, rectangle, minimum width=1.8cm, 
    minimum height=2cm, 
    text depth=2cm,
    text centered, 
    inner sep=1, 
    text width=2.0cm, font=\footnotesize,
    draw=black, 
    fill=blue!30},
    % https://tex.stackexchange.com/questions/6899/tikz-multi-line-text-in-the-node-description
    conNode/.style   = {draw, rectangle, minimum width=1.6cm, minimum height=1cm, text centered, text width=1.5cm, font=\footnotesize, draw=black,
    fill=yellow!30 
    % fill=white
    },
    dockerNode/.style = {draw, rectangle, minimum width=1.6cm, minimum height=0.5cm, text centered, text width=1.5cm, font=\footnotesize, draw=black, fill=yellow!30
    },
    cloudNode/.style = {draw,
    cloud, 
    cloud puffs=25.7, 
    cloud ignores aspect, 
    minimum width=3.5cm,
    minimum height=0.5cm, 
    align=center,
    fill=white
    },
}
\begin{figure}

\begin{subfigure}[b]{0.23\textwidth}
    %  https://tex.stackexchange.com/questions/115087/how-can-i-embed-an-external-image-within-a-tikzpicture-environment
    % https://tex.stackexchange.com/questions/91320/arrows-from-cloud-cross-the-outer-line
    % \begin{scaletikzpicturetowidth}{\textwidth}
    \begin{tikzpicture}[thick,scale=0.8, every node/.style={scale=0.8}]
    \coordinate (cloudco) at (1.6,1.5);
    \coordinate (serverAco) at (0,3.5);
    \coordinate (serverBco) at (2.5,3.5);
    \coordinate (clientco) at (0.8,0);
    % \draw[->] (cloudco) -- (90:2);
    % server node A
    \node [name=srvA, srv,draw] at (serverAco) {Edge Server A};
    \node [name=conA, conNode] at (srvA.center) {Offloading Container};
    \node [name=dockerA,dockerNode] at ([yshift=1em]srvA.south) {Docker};
    % draw line arrow from client to serverA
    \draw[->] (clientco) -- (dockerA.south);
    % server node B
    \node [name=srvB, srv] at (serverBco) {Edge Server B};
    \node [name=dockerB,dockerNode] at ([yshift=1em]srvB.south) {Docker};
    \node (phone) at (clientco)
        {\includegraphics[width=.2\textwidth]{figure/fig.jpg}};
    \node[cloudNode] (cloud) at (cloudco) {WAN};
    \draw[->] (phone) -- node[above]{moving} ++ (\textwidth/2,0);
    \end{tikzpicture} 
    \caption{Before Migration}
    \label{subfig:before}
\end{subfigure}
\begin{tikzpicture}[thick]
\draw[dashed] (0,0) -- (0,5.3);
\end{tikzpicture} 
\begin{subfigure}[b]{0.23\textwidth}
%  https://tex.stackexchange.com/questions/115087/how-can-i-embed-an-external-image-within-a-tikzpicture-environment
    % https://tex.stackexchange.com/questions/91320/arrows-from-cloud-cross-the-outer-line
    \begin{tikzpicture}[thick,scale=0.8, every node/.style={scale=0.8}]
    \coordinate (cloudco) at (1.6,1.5);
    \coordinate (serverAco) at (0,3.5);
    \coordinate (serverBco) at (2.5,3.5);
    \coordinate (clientco) at (1.5,0);
    % \draw[->] (cloudco) -- (90:2);
    % server node A
    \node [name=srvA, srv,draw] at (serverAco) {Edge Server A};
    \node [name=dockerA,dockerNode] at ([yshift=1em]srvA.south) {Docker};
    % server node B
    \node [name=srvB, srv] at (serverBco) {Edge Server B};
    \node [name=conB, conNode] at (srvB.center) {Offloading Container};
    \node [name=dockerB,dockerNode] at ([yshift=1em]srvB.south) {Docker};
    % draw line arrow from client to serverA
    \draw[->] (clientco) -- (dockerB.south);
    \node (phone) at (clientco)
        {\includegraphics[width=.2\textwidth]{figure/fig.jpg}};
    \node[cloudNode] (cloud) at (cloudco) {WAN};
    \draw[->] (phone) -- node[above]{moving} ++ (\textwidth/2,0);
    \end{tikzpicture}

\caption{After Migration}
\label{subfig:after}
\end{subfigure}
\caption{ Offloading Serivce Handoff: Before and After Migration of Offloading Container. }
\label{fig:before-after}
\end{figure}

In this section, we introduce the design of our service hand-off framework based on live migration of Docker containers. We will first briefly introduce a simple usage scenario and give an overview of system architecture and the algorithm of service handoff. Then in section \ref{design:mapImage}, we will discuss our methodology of storage synchronization based on Docker image layers' sharing between two different Docker hosts. Finally, in section \ref{design:memDiff}, \ref{design:compression}, and \ref{design:pipe}, we will discuss more strategies to further improve the migration speed by memory difference transfer, file compression, pipelined and paralleled processing during the migration of Docker containers.

\subsection{System Overview}

Figure~\ref{fig:before-after} shows an example usage scenario of offloading service hand-off based on Docker container migration. 
For example, the mobile client runs a real time face recognition application OpenFace\cite{openface2016} where it continuously sending the edge server images from the device's camera and the offloading containers processing the image via deep neural networks algorithm and send the name of the user back to the client as feedback.

Before migration, mobile client offloads its computations to its nearest edge server \textbf{A}, where all the computations are processed inside a Docker container. 
When the mobile client moves far away from server \textbf{A} and reaches the service area of edge server \textbf{B}, it's offloading computation will be migrated from server \textbf{A} to server \textbf{B}. This is done via migration of its offloading Docker containers, where all the runtime memory states as well as its storage data are synchronized to the target server \textbf{B}.


% https://tex.stackexchange.com/questions/350729/software-architecture-picture-in-tikz
\tikzset{
    srvos/.style  = {draw, rectangle, minimum width=3.5cm, 
    minimum height=0.75cm, 
    text width=2.9cm,
    text centered, 
    inner sep=1, 
    text width=2.0cm,
    % font=\footnotesize, 
    draw=black, 
    % fill=blue!30
    },
    dockerNode/.style = {draw, rectangle, minimum width=3.5cm, 
    minimum height=0.75cm,
    text width=2.9cm,
    % text depth=0.15cm,
    text centered, 
    % font=\footnotesize, 
    draw=black,
    % fill=yellow!30
    },
    % https://tex.stackexchange.com/questions/159440/tikz-cylinder-alignment-and-fit
    aufs/.style  = {draw, 
    cylinder, shape border rotate=90, 
    shape aspect=0.2,
    minimum height=1.1cm, 
    minimum width=1.4cm,
    text depth=0.7cm,
    % text height=1.5cm,
    % font=\footnotesize,
    draw=black, 
    % fill=blue!30
    },
    aufsImage/.style  = {draw, 
    cylinder, shape border rotate=90, 
    shape aspect=0.2,
    minimum height=0.5cm, 
    minimum width=0.5cm,
    % text depth=1cm,
    % text height=1.5cm,
    % font=\footnotesize,
    draw=black, 
    % fill=blue!30
    },
    % https://tex.stackexchange.com/questions/6899/tikz-multi-line-text-in-the-node-description
    conNode/.style = {draw, rectangle, minimum width=1.4cm, minimum height=2cm,
    text depth=1.5cm,
    % text height=1.5cm,
    text centered, text width=1.3cm,
    % font=\footnotesize, 
    draw=black,
    % fill=yellow!30 
    % fill=white
    },
    cloudNode/.style = {draw,
    cloud, 
    cloud puffs=18.7, 
    cloud ignores aspect, 
    minimum width=1.7cm,
    minimum height=1cm, 
    % font=\footnotesize, 
    align=center,
    fill=white
    },
    bigcloudNode/.style = {draw,
    cloud, 
    cloud puffs=25.7, 
    cloud ignores aspect, 
    minimum width=2cm,
    text width=2cm,
    minimum height=1cm, 
    align=center,
    fill=white
    },
    otherCloudNode/.style = {draw,
    cloud, 
    cloud puffs=25.7, 
    cloud ignores aspect, 
    minimum width=2cm,
    text width=2cm,
    minimum height=1cm, 
    align=center,
    fill=white
    },
}
\begin{figure*}
%  https://tex.stackexchange.com/questions/115087/how-can-i-embed-an-external-image-within-a-tikzpicture-environment
% https://tex.stackexchange.com/questions/91320/arrows-from-cloud-cross-the-outer-line
\begin{tikzpicture}[thick, 
    font=\footnotesize,
% scale=0.8, every node/.style={scale=0.8}
]
\coordinate (cloudco) at (3.5,1.5);
\coordinate (dockerHubco) at (8.5,6.5);
\coordinate (serverOtherco) at (11,4);
\coordinate (serverAco) at (0,1);
\coordinate (serverBco) at (7,1);
\coordinate (clientAco) at (0,0);
\coordinate (clientBco) at (7,0);
% \coordinate (clientco) at (0.8,0);
% \draw[->] (cloudco) -- (90:2);
% 
% Server OS Node A
\node [name=srvosA, srvos,draw] at (serverAco) {OS(Linux)};
% Docker Node
\node [name=dockerA,dockerNode] at ([yshift=0.9em]srvosA.north) {Docker Daemon};
% ConNode
\node [name=conA, conNode, anchor=south, draw] at ([yshift=2em, xshift=-2.5em]dockerA.north) {Offloading Container};
% AUFS
% \node [name=aufsA, above = 1cm of srvosA, aufs,draw] {AUFS};
\node [name=aufsA, aufs,anchor=south,draw] at ([yshift=2em, xshift=2.8em]dockerA.north) {AUFS};
% AUFS image
% https://tex.stackexchange.com/questions/50934/how-to-mark-a-node-in-a-graph-with-a-pattern
\node[pattern=bricks, aufsImage] at ([xshift=-1em, yshift=1em]aufsA.south) {};
\node[pattern=bricks, aufsImage] at ([xshift=1em, yshift=1em]aufsA.south) {};


% Server OS Node B
\node [name=srvosB, srvos,draw] at (serverBco) {OS(Linux)};
% Docker Node
\node [name=dockerB, dockerNode] at ([yshift=0.9em]srvosB.north) {Docker Daemon};
% AUFS
\node [name=aufsB, aufs, anchor=south, draw] at ([yshift=2em, xshift=-2.8em]dockerB.north) {AUFS};
% AUFS image
% https://tex.stackexchange.com/questions/50934/how-to-mark-a-node-in-a-graph-with-a-pattern
\node[pattern=bricks, aufsImage] at ([xshift=-1em, yshift=1em]aufsB.south) {};
\node[pattern=bricks, aufsImage] at ([xshift=1em, yshift=1em]aufsB.south) {};

% ConNode
\node [name=conB, conNode,anchor=south, draw] at ([yshift=2em, xshift=2.5em]dockerB.north) {Offloading Container};

 
% % DockerHub
% \node[ bigcloudNode] (dockerhub) at (dockerHubco) {Docker Hub};
% % notes under Docker hub
% \node[anchor=center, below = of dockerhub] {
% \ref{predownload}. Pre-Download Images
% };

 
% Other servers except B
 \node[otherCloudNode] (serverOther) at (serverOtherco) {Other Near Edge Servers};

% transparent client node A B
\node [name=clientA] at ([yshift=1em]clientAco){Source};
\node [name=clientB] at ([yshift=1em]clientBco){Target};


% pre download images
% https://tex.stackexchange.com/questions/99193/rotate-align-text-on-arrows-in-tikz
% https://tex.stackexchange.com/questions/27222/draw-an-arc-between-2-nodes-and-label-it-in-tikz
\draw[
% green, 
% dashed, 
bend right,
->] (dockerHubco) to
node[black, 
% sloped, 
anchor=center, above] {
% \ref{predownload}. Download Images
}
(aufsB.north);

\draw[
% green, dashed, 
bend left, ->] (dockerHubco) to 
node[black, 
% sloped,
anchor=center, above] {
% \ref{predownload}. Download Images
}
(serverOther.north);

 
% DockerHub
\node[ bigcloudNode] (dockerhub) at (dockerHubco) {Docker Hub};
% notes under Docker hub
\node[anchor=center, below = of dockerhub, text width = 3cm, yshift=2em, xshift=1em] {
\ref{predownload}. Pre-Download Images
};


% send pre dump mem
\draw[
% green, 
dashed, ->] 
([yshift=2.3em]conA.east) 
-- node[black, above
% , text width=2cm
]{
\ref{predump}. Send Pre-Dump Memory
} 
node[below, black
% , text width=2cm
]{\ref{img-sync}.
Memory Difference
}
([yshift=2.3em]conB.west);

% pre dump mem on A
\draw[<-] 
([xshift=-2.5em]dockerA.north) 
-- 
node[black, left,text width=2cm,align=left
% , text width=2cm
]{
\ref{predump}.Pre-dump
\\
\ref{checkpoint}. Final Dump
} 
([xshift=0em]conA.south);



% request migration
\draw[->] ([xshift=2em]clientAco) -- 
node[right, black]{
\ref{prepare}. Migration Request
}
([xshift=2em]srvosA.south);



% prepare
\draw[
% green, 
dashed, ->] ([yshift=0em]aufsA.east) -- 
node[above, black]{
\ref{prepare}. Match Image Layers
}
node[below, black]{
\ref{fs-sync}. Container Layer Sync
}
([yshift=0em]aufsB.west);

% ID remapping on source
\draw[<->] 
([xshift=2.8em]dockerA.north) 
-- 
node[black, right,text width=2cm,align=left
% , text width=2cm
]{
\ref{prepare}. ID Remapping
} 
([xshift=0em]aufsA.south);

% ID remapping & reload daemon on target
\draw[<->] 
([xshift=-2.8em]dockerB.north) 
-- 
node[black, left,text width=2cm,align=left
% , text width=2cm
]{
\ref{prepare}. ID Remapping
\\
\ref{daemon-Reload}. Reload Daemon
} 
([xshift=0em]aufsB.south);


% Memory Difference apply & restore.

\draw[<->] 
([xshift=1em]dockerB.north) 
-- 
node[black, right,
text width=4cm,
align=left
% , text width=2cm
]{
\ref{img-sync}. Apply Memory Difference
\\
\ref{restore}. ID Remapping
} 
([xshift=-1.5em]conB.south);

% WAN between A and B
 \node[cloudNode] (cloud) at (cloudco) {WAN};
 
\end{tikzpicture} 
\caption{System Architecture for Offloading Service Hand-off 
% Based on Docker Container's Migration. 
}
\label{fig:migstages}
\end{figure*}


Figure~\ref{fig:migstages} shows the design details of architecture as well as the migration algorithm from different processing stages described below:

\begin{enumerate}[label=\textbf{S\arabic*}]

\item 
% \label{prepare} 
\label{prepare} 
\textbf{
% Pre-download Application Image.
Synchronize Base Image Layers
} Upon the container starts and offers offloading service to the client, we also dynamically predict the possible target edge servers that are near to the client. Then we request those target servers to synchronize the base image layers for that container. 

\item \label{predump} \textbf{Pre-dump Container.} Before we received the migration request, we will try to dump a snapshot of the container's runtime memory and send to the possible target edge servers we found in stage \ref{prepare}. The container is still running during this stage. 

\item \label{request}\textbf{Migration Request Received on Source Server}. Once the migration request is initiated, the request will be sent to the source server, and we regard as the start point of the service handoff.
% \item \label{prepare} \textbf{Remapping Layer Cache IDs.} Upon receiving the hand-off request to a certain target server, prepare for file system synchronization by remapping the image layers' local identifications on the target server.

% \item \label{firstsync} Upon receiving migration request, compress the thin container layer and send to target.

\item \label{checkpoint} \textbf{Checkpoint and Stop Container.} Upon receiving the migration request, the source edge server will checkpoint and stop the container in order to send all the latest run-time states to the target edge server.
From this point the container will stop running on the source server and waiting to be restored on the target server.

\item \label{fs-sync} \textbf{Contaienr Layer Synchronization} After checkpointed the container, the file system will not be changed. In this step, we will send the container layer's contents to the target server. At the same time, all the checkpointed runtime states and configuration files, such as \textit{state.json, config.v2.json, etc.} are also transferred to the target server. 

\item \label{daemon-Reload} \textbf{Docker Daemon Reload} After we send all the runtime state and configuration files to the target node, Docker daemon on the target node still cannot recognize the container immediately.  This is because those runtime state and configuration files are created by another Docker daemon on source node, the daemon on target node doesn't have those configurations loaded into its runtime database. So we need to reload the Docker daemon in order to reload those runtime state and configuration files it just received from the source node.

\item \label{img-sync} \textbf{Get, Send, and Apply Memory Difference.} After we get checkpointed images from the final dump of the container, we then compare this final dump memory to the pre-dumped memory in stage \ref{predump}. We then could get the memory difference and send only these difference to the target server. 
% \item \label{finalsync} Synchronize the file system for the last time, updating only the changed files of the container layer before last synchronization in \ref{firstsync}.

\item \label{restore} \textbf{Restore Container.} Finally, we will restore container on the target host with all its runtime status retrieved. At the same time, the system will go to stage \ref{prepare} to prepare the next iteration of service handoff in the future.

\end{enumerate}


% (Before migration, both the source and target edge server have the application base images downloaded.)

% \subsection{Design Assumptions}
% 
% 

\subsection{Strategy to Synchronize Storage Layers} \label{design:mapImage}

As we introduced above, Docker's storage driver supports layered images where each layer represents file system differences. A running container's layered storage is composed by one writable container layer and several read only base image layers. 

The thin writable container layer stores all the files created or modified by the new created container from the base image. As long as the container is running, this layer is subject to change. So we postpone the synchronization of container layer after the container is stopped.

All the base image layers are read only inside container. They will not be changed during the whole life cycle of the container. So we could synchronize those layers as early as possible.

There are two kinds of base image layers. Firstly, in most cases, the base image layers for the container are downloaded by `docker pull` command from the centralized Registry like Docker Hub. All those images can be shared by downloading from the same Registry. Secondly, the container itself can also create its own image layers by saving its current container layer as one read only image layer. 

For these two kinds of base images, we use different strategies to synchronize. Generally, we will download the common image layers between two Docker hosts from the centralized image Registry and transfer only the different image layers between two Docker hosts.

By downloading the common image layer from the Registry, we could reduce the traffic between two edge servers. Furthermore, the download could start as soon as the container is created on the source server, and thus the download time could be masked ahead of migration starts.

Finally, for the base image layers create by the container, we transfer each image layer upon the image layer is created, no matter whether the migration has started or not. 

\subsection{Layer ID Remapping} \label{idremapping}

As discussed in section \ref{intro:idMatching} and \ref{intro:aufs:layerIDMapping}. The downloaded images from the common Registry have different cache IDs exposed to each Docker host. In order to share these common images across different Docker host, we need to match these  image layers by its original IDs instead of the cache IDs.

In order to do this, we first remapping the cached IDs to its original IDs and compare the original IDs on two Different Docker hosts. If the two image layers on two host share the same original ID, we could determine they are exactly the same image layer. 

Then, for the matched original layer IDs on both of the Docker hosts, we will to remap the original IDs to its local cache IDs on the target host. Now that we get the new cache IDs on the target Docker host. Then we replace the old cache IDs on the source host with the new cache IDs on the target Docker host.  

By doing so, the image layer stack list for the migrated container will be reset with the new cache IDs that is addressable for the Docker host on the target server. So that when we restore the container in the future, its file system will be mounted correctly from the shared image layers on the target server.


% since the target Docker daemon address each layer with its locally generated cache IDs instead of the original IDS. 

% in order to make it compatible with the local addressing system of that Docker daemon. 
% After remapping all shared image layers to their local cache IDs on the target Docker, we need also to reset the image layer stack list for the migrated container we sent to the target, so that the container's file system will be mounted correctly on the target Docker host.

% The base image layer could either be downloaded from 
% the prepare for file system synchronization by remapping the image layers' local identifications on the target server.
% download the Docker image of the offloading application from Docker Hub.


\subsection{Pre-Dump \& Dirty Memory Synchronization} \label{design:memDiff}

In order to reduce the total memory image size during hand-off, we checkpoint the container and dump a snapshot of the container's memory in stage \ref{predump}. This could happen as soon as the container is created, or more enlightened, we could dump the memory when the most frequently used binary programs in the application are loaded into the memory. This snapshot of memory will be served as the base memory image for the live migration that happens next.

After the base memory image is dumped, it is transferred immediately to the server. We assume the transfer will be finished before the hand-off starts. It's reasonable since we can send the base memory image as early as seconds after the container starts. 
Also, the application's Docker images on the near edge servers are also started to download after the container starts and before the hand-off starts, we could process those two steps in parallel to reduce the total time cost. This is further discussed in section \ref{design:pipe}. Therefore, upon the hand-off starts, we can have the base memory image of the container already on the target server. 

Then after the client being served at the source node and starts to move into another edge server's service area, a migration request will be sent to the source node. 
% As for who send the migration request, it's out of the scope of this paper.

Once the migration request is received, we will  checkpoint and stop the container. After get the checkpointed memory, we do a diff operation on it from the base memory image we dumped in stage \ref{predump}. 
Then we only need to transfer the difference of the dump memory to the target node. 

\begin{figure}[b]
\begin{subfigure}[b]{0.23\textwidth}
    % \centering
    \includegraphics[width=\linewidth]
    % {figure/test-tarssh.eps}
    {figure/dumps/face-518744dump_iter_11.eps}
    \caption{Memory Snapshot Size}
    \label{fig:dump}
\end{subfigure}
\begin{subfigure}[b]{0.23\textwidth}
    % \centering
    \includegraphics[width=\linewidth]
    % {figure/test-tarssh.eps}
    {figure/dumps/face-518744-iter-11.eps}
    \caption{Dirty Memory, original}
    \label{fig:diff-ori}
\end{subfigure}
\begin{subfigure}[b]{0.23\textwidth}
    % \centering
    \includegraphics[width=\linewidth]
    % {figure/test-tarssh.eps}
    {figure/dumps/face-518744-iter-11-2.eps}
    \caption{Dirty Memory, adjacent}
    \label{fig:diff-adj}
\end{subfigure}

\begin{subfigure}[b]{0.23\textwidth}
    % \centering
    \includegraphics[width=\linewidth]
    % {figure/test-tarssh.eps}
    {figure/dumps/face-518744diffdiff_ori11.eps}
    \caption{Dirty Difference, original}
    \label{fig:diffdiff-adj}
\end{subfigure}
\begin{subfigure}[b]{0.23\textwidth}
    % \centering
    \includegraphics[width=\linewidth]
    % {figure/test-tarssh.eps}
    {figure/dumps/face-518744diffdiff_adj11.eps}
    \caption{Dirty Difference, adjacent}
    \label{fig:diffdiff-ori}
\end{subfigure}

\caption{ Dirty Memory Size Analysis for OpenFace. 
The data is collected from 11 dumps of the running container of OpenFace, labeled from dump 0 to 10. Memory is dumped every 10 seconds. Container continues to run during the first 10 dumps, and stops after 11th dump.
Figure \ref{fig:dump} shows checkpointed memory size for total 11 dumps.
% Container is checkpointed and restored immediately in order to simulate the --leave-running options of checkpoint command.
Figure \ref{fig:diff-ori} shows dirty memory size beween each of dump 1~10 and the original dump 0.
Figure \ref{fig:diff-adj} shows dirty memory size between two adjacent dumps.
Figure \label{fig:diffdiff-adj} shows the difference of dirty memories between three adjacent dumps.
Figure \ref{fig:diffdiff-ori} shows the difference of dirty memories compared to the original dirty memory from dump 0 to dump 1. } 
\end{figure}



\subsection{Data Transfer}\label{design:compression}

% During migration, the files need to be transferred includes both text files as well as binary files. 
During the live migration of a container, we mainly have 4 kinds of data need to transfer:
the layer stacks information, the file system of thin writable \textit{container layer}, the meta data files of the container, and snapshot of the container's memory and image difference. Some of the data is in the form of messages, such as layer stacks information. Some of them are plain text files, such as most contents. And the memory snapshots and image differences are all binary image files. According to the feature of the files, we use different strategies to transfer the data.

The layer stacks information is send via UNIX RPC API implementation in \cite{phaul}, which is based on a socket connection created by python scripts. Those information is only a list of SHA256 ID strings, so it's quite efficient to be sent as a socket message. There is no need to compress it since the overhead of compression would mask its benefits.

As for other data, including the container writable layer, meta data files, the dump memory images, and image differences, we use bzip2 to compress and send via authorized ssh connection.

% \ref{predump}, \ref{img-sync}.

% For the root file system and meta data files. We have two stages of files synchronization before and after the container was checkpoints. The first synchronization in \ref{firstsync} will transfer the base file system for the container layer, so it's better to compress it. We use 'tar' command to compress and sent via SSH. For the final synchronization of the file system, since we already have most of the files transferred, we choose to use 'rsync' to send over only the changed files.

% For the checkpointed memory images, we compress them via 'tar' compression and send via SSH connections.

% Syncing smaller files individually through rsync or scp results in each file starting at least one own data packet over the net. If the file is small and the packets are many, this results in increased protocol overhead. Now count in that there are more than one data packets for each file by means of rsync protocol as well (transferring checksums, comparing...), the protocol overhead quickly builds up

% \url{http://unix.stackexchange.com/questions/30953/tar-rsync-untar-any-speed-benefit-over-just-rsync}

% TODO: more details about compression and synchronization: more memory images compression techniques (to replace tar cmd). 

% TODO: rsync vs. tar comparison


\subsection{Parallel \&  Pipelined Processing}\label{design:pipe}

With the help of parallel and pipelined processing,  we could improve the efficiency of the whole process and reduce the total migration time.

% From stage \ref{prepare} to stage \ref{restore}, there are lots of tasks that could be run in parallel between the source and target servers.

Firstly, as we mentioned above in \ref{design:memDiff}, starting a container will trigger two events to happen in parallel: downloading/synchroizing images from Registry in cloud, and pre-dumping/sending base memory images to the possible edge servers near the source node. Those two process could be run at the same time in order to reduce the total time of stage \ref{prepare} and \ref{predump}. When the possible target servers are downloading images from Registry, the source server will checkpointing and sending the running container's base memory images to the possible target server.

Secondly, daemon reload in stage \ref{daemon-Reload} is happend on target host, it could be processed while the source host is sending the memory difference to the target host in stage \ref{img-sync}. Stage \ref{fs-sync} cannot be paralleled with \ref{daemon-Reload}, because daemon reload on the target host requires the configuration data files sent in stage \ref{fs-sync}.

Thirdly,
in stage \ref{fs-sync}, we use compression to send all files in the \textit{container layer} over an authorized ssh connection between the source and target host. The compression and transfer of the \textit{container layer} can be pipelined using Linux pipes.

Lastly, in stage \ref{img-sync}, after we get the final memory snapshot, we will need to get memory difference by comparing the base memory images with the images in the new snapshot, then we send the difference to target and patch the difference to the base memory image on the target host. This whole process could also be piplined using Linux pipes. 


\section{Evaluation}


In this section, we evaluate the performance of our system under different workloads and different network conditions. 
Specifically, we evaluate offloading with our service handoff system by measuring the degree to which total hand-off time and total transferred size are reduced.

\subsection{Set-Up and Benchmark Workloads}

Our test bed is built on a desktop server, which is equipped with Intel Core i3-6100 Processor (3.70GHz, 2 cores, 4 threads), and 16GB DDR4 memory. Two virtual machines are running with 2 vcpus and 4GB memory each.
% The laptop is with Intel$^{R}$ Core$^{TM}$ Duo T6570 (2.2GHz, 2 cores) and 2GB DDR2 memory. Both the machines and virtual machines are running Ubuntu 16.04 LTS as the operating system.
A laptop acts as a client, this has Intel Core i5-2450M CPU (2.5GHz, 2 cores, 4 threads) and 4GB DDR3 memory. 
The hosts and virtual machines are running Ubuntu 16.04 LTS as the operating system. Docker version is 1.10-dev which is built with experimental feature enabled.

In our experiment, we set up migration scenarios using 
emulated bandwidth on two virtual machines each running a Docker instance.
Docker containers are migrated from the Docker host on source VM to the Docker host on another VM. 
Two virtual machines run on the same physical server. 

We use Linux Traffic Control (tc\cite{tc} )  tool to control network traffic. In order to test our system under WANs , we emulated low bandwidth ranging from 5Mbps to 45Mbps. According to the average bandwidth observed on the Internet\cite{internet2014state}, we set  latency at a fixed 50ms to emulate the WAN environment for edge computing.

Since edge computing environments can also be adapted to a relatively large LAN network such as a university, or other large community, it usually can reach higher bandwidth as well as lower latency. Therefore, we also tested several higher bandwidths, ranging from 50Mpbs to 500Mpbs. The latency is set to 6ms which is an average latency from the author's university LAN.

For the offloading workloads, we use Busybox as a simple workload to show the functionality of the system as well as the non-avoidable system overhead when processing container migration. In order to show the realtime performance of offloading service handoff, we chose OpenFace to provides a sample workload.

% \subsubsection{Busybox}



% \subsubsection{OpenFace
% % : A Realtime Wearable Cognitive Application
% }

\subsection{Dirty Memory \& Experimental Findings}

\begin{figure}[b]
\begin{subfigure}[b]{0.3\textwidth}
    % \centering
    \includegraphics[width=\linewidth]
    % {figure/test-tarssh.eps}
    {figure/dumps-busybox/busybox-d34455dump_iter_11.eps}
    \caption{Memory Snapshot Size}
    \label{fig:busydump}
\end{subfigure}

\begin{subfigure}[b]{0.23\textwidth}
    % \centering
    \includegraphics[width=\linewidth]
    % {figure/test-tarssh.eps}
    {figure/dumps-busybox/busybox-d34455-iter-11.eps}
    \caption{Dirty Memory, original}
    \label{fig:busydiff-ori}
\end{subfigure}
\begin{subfigure}[b]{0.23\textwidth}
    % \centering
    \includegraphics[width=\linewidth]
    % {figure/test-tarssh.eps}
    {figure/dumps-busybox/busybox-d34455-iter-11-2.eps}
    \caption{Dirty Memory, adjacent}
    \label{fig:busydiff-adj}
\end{subfigure}
% \begin{subfigure}[b]{0.23\textwidth}
%     % \centering
%     \includegraphics[width=\linewidth]
%     % {figure/test-tarssh.eps}
%     {figure/dumps-busybox/busybox-d34455diffdiff_ori11.eps}
%     \caption{Dirty Difference, original}
%     \label{fig:busydiffdiff-adj}
% \end{subfigure}
% \begin{subfigure}[b]{0.23\textwidth}
%     % \centering
%     \includegraphics[width=\linewidth]
%     % {figure/test-tarssh.eps}
%     {figure/dumps-busybox/busybox-d34455diffdiff_adj11.eps}
%     \caption{Dirty Difference, adjacent}
%     \label{fig:busydiffdiff-ori}
% \end{subfigure}

\caption{ Dirty Memory Size Analysis for Busybox. 
Figure \ref{fig:busydump} shows checkpointed memory size for total 11 dumps.
% Container is checkpointed and restored immediately in order to simulate the --leave-running options of checkpoint command.
Figure \ref{fig:busydiff-ori} shows dirty memory size beween each of dump 1~10 and the original dump 0.
Figure \ref{fig:busydiff-adj} shows dirty memory size between two adjacent dumps.
% Figure \label{fig:busydiffdiff-adj} shows the difference of dirty memories between three adjacent dumps.
% Figure \ref{fig:busydiffdiff-ori} shows the difference of dirty memories compared to the original dirty memory from dump 0 to dump 1. 
} 
\label{fig:busydiff}
\end{figure}



\begin{figure}[b]
\begin{subfigure}[b]{0.23\textwidth}
    % \centering
    \includegraphics[width=\linewidth]
    % {figure/test-tarssh.eps}
    {figure/dumps/face-518744dump_iter_11.eps}
    \caption{Memory Snapshot Size}
    \label{fig:dump}
\end{subfigure}
\begin{subfigure}[b]{0.23\textwidth}
    % \centering
    \includegraphics[width=\linewidth]
    % {figure/test-tarssh.eps}
    {figure/dumps/face-518744-iter-11.eps}
    \caption{Dirty Memory, original}
    \label{fig:diff-ori}
\end{subfigure}
\begin{subfigure}[b]{0.23\textwidth}
    % \centering
    \includegraphics[width=\linewidth]
    % {figure/test-tarssh.eps}
    {figure/dumps/face-518744-iter-11-2.eps}
    \caption{Dirty Memory, adjacent}
    \label{fig:diff-adj}
\end{subfigure}

\begin{subfigure}[b]{0.23\textwidth}
    % \centering
    \includegraphics[width=\linewidth]
    % {figure/test-tarssh.eps}
    {figure/dumps/face-518744diffdiff_ori11.eps}
    \caption{Dirty Difference, original}
    \label{fig:diffdiff-adj}
\end{subfigure}
\begin{subfigure}[b]{0.23\textwidth}
    % \centering
    \includegraphics[width=\linewidth]
    % {figure/test-tarssh.eps}
    {figure/dumps/face-518744diffdiff_adj11.eps}
    \caption{Dirty Difference, adjacent}
    \label{fig:diffdiff-ori}
\end{subfigure}

\caption{ Dirty Memory Size Analysis for OpenFace. 
The data is collected from 11 dumps of the running container of OpenFace, labeled from dump 0 to 10. Memory is dumped every 10 seconds. Container continues to run during the first 10 dumps, and stops after 11th dump.
Figure \ref{fig:dump} shows checkpointed memory size for total 11 dumps.
% Container is checkpointed and restored immediately in order to simulate the --leave-running options of checkpoint command.
Figure \ref{fig:diff-ori} shows dirty memory size beween each of dump 1~10 and the original dump 0.
Figure \ref{fig:diff-adj} shows dirty memory size between two adjacent dumps.
Figure \label{fig:diffdiff-adj} shows the difference of dirty memories between three adjacent dumps.
Figure \ref{fig:diffdiff-ori} shows the difference of dirty memories compared to the original dirty memory from dump 0 to dump 1. } 
\end{figure}



Figure \ref{fig:facediff}
and  \ref{fig:busydiff} 
give an overview of the dump memory image sizes while running the container of OpenFace and Busybox. 
The data is collected from 11 dumps of the running container of Busybox and OpenFace, labeled from dump 0 to 10. Memory is dumped every 10 seconds. Container continues to run during the first 10 dumps, and stops after 11th dump. From the figure we can see the memory difference is much smaller than the original memory dump. The migration could be sped up by only  transferring the memory difference. 

Furthermore, 
Figure  
\ref{fig:busydiff-ori} $\sim$
\ref{fig:busydiff-adj}  
% \ref{fig:busydiffdiff-ori}  
and 
\ref{fig:facediff-ori} $\sim$
% \ref{fig:facediffdiff-ori} 
\ref{fig:facediff-adj} 
show us the sizes of memory differences between adjacent dumps as well as the original dump. We find a feature of memory access patterns for Busybox and OpenFace application. Although their memory is continuously changing, the changes reside in a specified area: a 4KB area for Busybox and 25MB area for OpenFace. 
In this case, iterative transferring of memory difference to the target server will be waste of effort. Although this will not happen on every application, it is valuable to be considered as a special case for further research regarding live migration of containers by iterative transfer. One way to utilize those kinds of memory access patterns is by not iteratively transferring every memory difference. Instead, transfer only the last iterative difference to achieve the same end result.

\subsection{Evaluation of Pipeline Performance}


\begin{figure}[tb!]
    \centering
    \includegraphics[width=0.48\textwidth]
    % {figure/test-tarssh.eps}
    {figure/evaluations/compare-pipe-nopipe-busybox.eps}
    \caption{Busybox: Time of Container Migration Stages with and without Pipelined Processing. }
    \label{fig:busypipe}
\end{figure}

\begin{figure}[tb!]
    \centering
    \includegraphics[width=0.48\textwidth]
    % {figure/test-tarssh.eps}
    {figure/evaluations/compare-pipe-nopipe-face.eps}
    
    \caption{OpenFace: Time of Container Migration Stages with and without Pipelined Processing
    % Pipelined Processing of ImgDiff, ImgSend, and MemApply 
    }
    \label{fig:facepipe}
\end{figure}

In order to verify the effectiveness of pipelined processing, we implemented the pipeline processing of two time consuming steps: \textit{imgDiff} and \textit{imgSend}, where \textit{imgDiff} is to get memory difference and \textit{imgSend} is to send memory difference to the target during migration. 
Figure \ref{fig:facepipe} and Figure \ref{fig:busypipe} reports the timing benefits we could achieve by using pipelined processing. From the figure, we can see that, without pipelined processing, most time costs are by getting and sending the memory difference. After applying pipelined processing, we could save $5\sim 8$ seconds for OpenFace migration. Busybox also saves a certain amount of time with pipelined processing.


\subsection{Overall Performance and Comparison with State-of-the-Art}
We evaluate the total \textit{handoff time} and \textit{transferred size } during offloading service handoff under different bandwidths and network latencies. 

Table \ref{tab:all} shows an overview of the performance of our system under different network bandwidth conditions.
% 
\begin{figure}[tb!]
    \centering
\begin{subfigure}[b]{0.23\textwidth}
    % \centering
    \includegraphics[width=\linewidth]
    % {figure/test-tarssh.eps}
    {figure/eva-individual/predump-docker-pipe-busybox-stack-bar.eps}
    \caption{Busybox}
    \label{fig:allband-busy}
\end{subfigure}
\begin{subfigure}[b]{0.23\textwidth}
    % \centering
    \includegraphics[width=\linewidth]
    % {figure/test-tarssh.eps}
    {figure/eva-individual/predump-docker-pipe-face-stack-bar.eps}
    \caption{OpenFace }
    \label{fig:allband-face}
\end{subfigure}
    \caption{Time of Container Migration Under Different Bandwidths (WANs \& LANs).
    }
    \label{fig:allband}
\end{figure}

% Figure \ref{fig:allband}.
% Figure \ref{fig:allband-busy} 
% Figure \ref{fig:allband-face}



TODO one large table here:

\begin{table}[h]
\centering
\begin{tabular}{|P{3.5em}|P{2em}|P{3em}|P{2.8em}|P{3.5em}|P{3.5em}|}
\hline
\multicolumn{2}{|P{5em}|}{Bandwidth (Mbps)} &  Handoff Time(s) & Down Time(s) & Pre-Transfer Size & Final-Transfer Size \\                \hline
\multirow{13}{*}{Busybox} 
 &5     &3.90  & 3.23  & & c\\
 &15    &X  &   & c &c \\
 &20    &x  &x  & c &c \\
 &25    &x  &x  & c &c \\
 &30    &x  &x  & c &c \\
 &35    &x  &x  & c &c \\
 &40    &x  &x  & c &c \\
 &45    &x  &x  & c &c \\
 &50    &x  &x  & c &c \\
 &100    &x  &x  & c &c \\
 &200    &x  &x  & c &c \\
 &500    &x  &x  & c &c \\
 &1000    &x  &x  & c &c \\\cline{1-6}
\multirow{13}{*}{Busybox}
 &5     &n  &   & c &c \\
 &15    &X  &   & c &c \\
 &20    &x  &x  & c &c \\
 &25    &x  &x  & c &c \\
 &30    &x  &x  & c &c \\
 &35    &x  &x  & c &c \\
 &40    &x  &x  & c &c \\
 &45    &x  &x  & c &c \\
 &50    &x  &x  & c &c \\
 &100    &x  &x  & c &c \\
 &200    &x  &x  & c &c \\
 &500    &x  &x  & c &c \\
 &1000    &x  &x  & c &c \\\cline{1-6}
\end{tabular}
    \caption{
        Average of 10 runs and relative standard deviations (RSDs, in parentheses) are reported. 
        Handoff time is from the time the server receives migration request until migration done; 
        The down time is the time from the container is checkpointed to the container is restored on the target.
    }
    \label{tab:perform}
\end{table}

--------------------------------------

bandwidth/app   migration time / down time  /  FS size / Image size

5Mbps, 50ms

 \quad busybox
    
 \quad   face
    
  \quad  ...
    
10Mbps, 50ms

 \quad busybox
    
 \quad   face

  \quad  ...

15Mbps, 50ms

 \quad busybox
    
 \quad   face

  \quad  ...

25Mbps, 50ms

 \quad busybox
    
 \quad   face

  \quad  ...


30Mbps, 50ms

==========================================

% \begin{table}[!t]
% \centering

% \caption{Optimized using both layered images and compression with 5 Mbps bandwidth and 50 ms latency}
% \label{table_samehost_opt}


% % Some packages, such as MDW tools, offer better commands for making tables
% % than the plain LaTeX2e tabular which is used here.
% % \begin{tabular}{|l|l|l|l|l|l|}
% \begin{tabular}{|M{1.2cm}|M{0.9cm}|M{0.9cm}|M{1.3cm}|M{1.1cm}|
% % M{1.0cm}|
% }
% \hline
% Container & {Total time}  & {Down time} & { FS size} & { Mem size } 
% % & { Total Size }
% \\ \hline
% Busybox & $4.13s$ & $4.10s$ & $14KB$ & $29KB$ \\\hline
% OpenFace & $22.14 s$ & $22.00s$ & $100KB$ & $102.1MB$ \\\hline
% \end{tabular}
% \end{table}

% \begin{table}[!t]
% % increase table row spacing, adjust to taste
% % \renewcommand{\arraystretch}{2.5}
% % if using array.sty, it might be a good idea to tweak the value of
% % \extrarowheight as needed to properly center the text within the cells
% \centering

% \caption{Optimized using both layered images and compression with 5 Mbps bandwidth and 50 ms latency}
% \label{table_ethernet_opt}

% % Some packages, such as MDW tools, offer better commands for making tables
% % than the plain LaTeX2e tabular which is used here.
% % \begin{tabular}{|l|l|l|l|l|l|}
% \begin{tabular}{|M{1.2cm}|M{0.9cm}|M{0.9cm}|M{1.3cm}|M{1.1cm}|
% % M{1.0cm}|
% }
% \hline
% Container & {Total time}  & {Down time} & { FS size} & { Mem size } 
% % & { Total Size }
% \\ \hline 
% Busybox & $5.24s$ & $4.4s$ & $14KB$ & $29KB$ \\\hline
% OpenFace & $29.93s$ & $?s$ & $101KB$ & $102.1MB$ \\\hline
% \end{tabular}

% \end{table}


% \begin{table}[!t]
% % increase table row spacing, adjust to taste
% % \renewcommand{\arraystretch}{2.5}
% % if using array.sty, it might be a good idea to tweak the value of
% % \extrarowheight as needed to properly center the text within the cells
% \centering

% \caption{Optimized using both layered images and compression with 5 Mbps bandwidth and 50 ms latency}
% \label{table_wireless_opt}

% % Some packages, such as MDW tools, offer better commands for making tables
% % than the plain LaTeX2e tabular which is used here.
% % \begin{tabular}{|l|l|l|l|l|l|}
% \begin{tabular}{|M{1.2cm}|M{0.9cm}|M{0.9cm}|M{1.3cm}|M{1.1cm}|
% % M{1.0cm}|
% }
% \hline
% Container & {Total time}  & {Down time} & { FS size} & { Mem size } 
% % & { Total Size }
% \\ \hline
% Busybox & $7.12s$ & $6.80s$ & $14KB$ & $29KB$ \\\hline
% OpenFace & $60.23s$ & $60.20s$ & $101KB$ & $102.1MB$\\\hline
% \end{tabular}

% \end{table}



\textit{Handoff time}
is from the time the source host receives a migration request until the offloading container is successfully restored on the target host. The \textit{Down time} is from the time when the container is stopped on the source server to the time when the container is restored on the target server. \textit{Handoff time} and \textit{Down time} are very similar since we immediately checkpoint and stop the container once we get the migration request. 

\textit{Pre-Transfer Size} is the transferred size before \textit{handoff} starts, i.e. from stage \ref{prepare} until stage \ref{request} . \textit{Final-Transfer Size} is the transferred size during the handoff, i.e. from stage \ref{request} until the end of final stage \ref{restore}.


\begin{figure}[tb!]
    \centering
    \includegraphics[width=0.48\textwidth]
    % {figure/test-tarssh.eps}
    {figure/evaluations/compare-50ms-6ms-busybox.eps}
    \caption{Busybox: Comparison of Migration Time Under Latency of 50ms and 6ms. }
    \label{fig:busylatency}
\end{figure}

\begin{figure}[tb!]
    \centering
    \includegraphics[width=0.48\textwidth]
    % {figure/test-tarssh.eps}
    {figure/evaluations/compare-50ms-6ms-face.eps}
    
    \caption{OpenFace: Comparison of Migration Time Under Latency of 50ms and 6ms. 
    }
    \label{fig:facelatency}
\end{figure}

Figure \ref{fig:busylatency}
and
Figure \ref{fig:facelatency} shows the performance under difference network latencies of $50ms$ and $6ms$ for Busybox and OpenFace.  It shows a tiny difference when facing different latencies. This means our system is suitable for a wide range of network latencies.


From Table \ref{tab:all} and Figure \ref{fig:busylatency} , we can see the simple Busybox container can be migrated very quickly regardless the network bandwidth and latency. 

From Table \ref{tab:all} and Figure \ref{fig:facelatency} , we can see the OpenFace offloading container can be migrated within $49$ seconds under the lowest bandwidth $5$Mbps with $50$ ms latency. Compared to the work in \cite{ha2015vmhandoff}, which has a handoff time of $247$ seconds under the same network conditions. For $25$Mbps and $50$ ms latency, our system achieves $17.4$ seconds while  \cite{ha2015vmhandoff}  needs $39$ seconds. The relative standard deviations in Table \ref{tab:all} shows the robustness of our experimental result.


Therefore, our system could reduce the total handoff time by $56\%\sim 80\%$ compared to the state-of-the-art work of VM handoff \cite{ha2015vmhandoff} on edge computing platforms. 






% \subsubsection{Performance Under Different Bandwidths}

% \subsubsection{Performance Under Different Latencies}







    % \item Laptop-Wireless-WAN:
    % The migration source host is a virtual machine on the Laptop, another is the virtual machine on the server as in 1). Both virtual machines are running the same version of Docker.  They are connected through WAN network with about 2 miles geological distance. Each virtual machines is running within its own LAN network. The virtual machine on Laptop connects to the internet through wireless adapter. The virtual machine on the server is connected to the Internet via Ethernet. 
% \end{enumerate}



% \section{Related Work}

% In this section, we discuss relate work in the areas of dynamic VM synthesis, VM handoff, and container migration.

\subsection{Dynamic VM Synthesis and VM Hand-off} \label{motivation:vmhandoff}

% Many researches have been done in order to accelerate the service handoff across edge servers.
Techniques based on virtual machine migration have been proposed in  \cite{ha2015vmhandoff} \cite{satya2009case} to accelerate the service handoff across edge servers. 
% As discussed in irtual Machine live migration in data centers could not be used for service handoff.
% which was well discussed in \cite{ha2015vmhandoff}. Firstly, Live migration of VM requires high bandwidth of connection and transfers Gigabytes of data in each migration, where service handoff for offloading happens on low bandwidth WANs. Secondly, live migration in data centers aims at short duration of the VM down time. However, service handoff across edge servers requires to optimize the time of total operations from the beginning of handoff request. 

In order to enable high speed service handoff based on virtual machine migration techniques, Satyanarayanan \textit{et al.} in \cite{satya2009case} proposed VM synthesis to divide a huge VM imges into a base VM image and a relatively small overlay image for one specific application. Before migration starts, the server is assumed to have a base image, so that the transition of this base image is avoided. During migration, only the application's overlay image is transferred from the source to the target server.

Based on the work of VM synthesis, VM handoff across Cloudlet servers (alias of edge servers) was proposed in \cite{ha2015vmhandoff}. VM handoff allows the mobile client to seamlessly transfer the runtime status of its VM from one edge sever to a nearer edge server. Besides VM synthesis, it uses dirty tracking for both file system and main memory  to reduce the transferred size during handoff.
Additionally, it also uses compression, pipelined processing to reduce transfer time and algorithms to adapt to different network conditions.
Finally, it highly reduced the transferred size and migration time under WAN network environments.

While significant efforts have been invested into VM synthesis and handoff, the performance cannot meet the requirements of offloading services in edge computing environments, given the mobility of clients, dominance of wireless access, geographical distribution of edge nodes, and real-time interaction requirement. Instead,  our work is along the line of container, a lightweight OS-level virtualization technique, and focus on the migration of container to reduce the total handoff time. 
% However,
% % the transfer size of the application still had tens or hundreds of megabytes on average. 
% the total handoff time was still several minutes on a WAN network. For example, it will requires $245.5$ seconds to migrate a running OpenFace instance under 5Mbps bandwidth (50ms latency) network in \cite{ha2015vmhandoff}. 

% One of the reasons for the long latency of handoff the large image size we need to transfer. Although the VM synthesis could reduce the image size by splitting images into multiply layers and only transfer the application-specific layer, the total transferred size is still in the magnitude of tens of megabytes and even hundreds of megabytes. 
% The application layer was encapsulated with the whole application, both its static binary programs and its runtime memory data, which we think is an unnecessary cost. 

% On the other hand, the deployment of the VM synthesis system is challenging for the legacy system. In order to enable VM synthesis, the hypervisor of VMs need to be patched to track the dirty memory at runtime. Also the storage of VM images need to be adapted to Linux FUSE interfaces in order to track file system changes inside the running VMs. Those two changes are hard to deploy in practice since they changes the behavior of hypervisors and file systems of the legacy platform and will also cause lots of performance overhead. 

% \subsection{Latency Sensitive Mobile Applications}

% \subsubsection{Enlightened VM Post-Copy} In order to improve the 
% \subsection{Linux Containers}

\subsection{Containers Migration}
% Useful link: https://coreos.com/rkt/docs/latest/rkt-vs-other-projects.html#rkt-vs-runc

% Containers provide operating system level virtualization by running processes in isolated environments, 
Containers provide operating system level virtualization by running a group of processes in isolated environments. It's supported by the kernel features of namespace 
 and cgroup (control groups)\cite{rosen2014container}.
Namespaces are used to provide a private view of system resources for processes \cite{pike1992plan9}, while cgroups are used to restrict how many resources a group of processes can be accessed\cite{rosen2013namespaceIntro}.  

%  every resource in the system, either local or remote, is represented by a hierarchical file system; and a user or process assembles a private view of the system by constructing a file name space that connects these resources.

Container runtime is a tool that provides easy-to-use API for managing containers by abstracting the low-level technical details of namespaces and cgroups. 
Although creating a container by crafting cgroup and namepaces step by step is possible \cite{conScratch}, managing containers by container runtime is much more easier.  Such tools include LXC/LXD\cite{LXC} ,
runC\cite{runc},
rkt\cite{rkt}, 
OpenVZ\cite{openvz}, 
Docker\cite{docker}, etc. Different container runtime has different scenerios of usage. For example, LXC/LXD only cares about full system containers and doesn't care about what kind of application runs inside the container, while Docker aims to encapsulate a specific application in the container.  

Migration of containers become possible when CRIU (Checkpoint/Restore In Userspace)\cite{criu} supports the checkpoint/restore functionality for Linux. 
Now CRIU supports the checkpoint and restore of containers for OpenVZ, LXC, and Docker. 


Based on CRIU, OpenVZ now supports migration of containers. The implementation can be found from the CRIU community's open source project, P.Haul(process Hauler)\cite{phaul}.  It's claimed that the migration could be done within 5 seconds\cite{livmigcon}.
However, OpenVZ uses distributed storage system\cite{openvzfs}, where all files are shared across high bandwidth network. 
This means when it migrates the containers, it only needs to migrate the checkpointed memory images and no file transfer need to be done. 
However, due to the limited WAN bandwidth for edge servers, it's not possible to deploy distributed storage. Therefore, migration of OpenVZ containers is not suitable for service handoff on edge computing platforms.

For LXC containers, the migration is implemented in Qiu's thesis work \cite{qiu2016evaluating}, which is also based on the CRIU project. However, LXC regards container as a whole system container, and they don't have layered storage for containers. Therefore, when containers are migrated, all the file system for that container must be migrated together will all its memory status. The total transfer size will be comparable with migrating a whole virtual machine(VM). So, it's still not feasible for edge computing environment. 

For Docker containers, there is a sample migration helper in P.Haul's source code, which supports only an older version of docker-1.9.0-dev. And it's also extended by Ross Boucher to support docker-1.10-dev. However, as we have discussed in section \ref{migpractice}, we find this implementation transmits the entire file system of that container, regardless of the layered storage of Docker platform, which makes the migration unsatisfactorily slow across the edges of WAN network. Therefore, in this paper, we dive into the inner details of Docker storage and find out how to leverage the layered storage of Docker platform during the migration and finally avoid the transfer of unnecessary file system by sharing Docker images between different Docker host. 

%  a lightweight universal runtime container. runC is a low-level container runtime and an implementation of the Open Container Initiative specification. runC exposes and expects a user to understand low-level details of the host operating system and configuration.



\section{Conclusion and Future Work}

We propose a framework that enhances service handoff across edge offloading servers by leveraging the layered storage system of Docker containers to improve migration performance.
Our system enables the edge computing platform to continuously provide offloading services with low end-to-end latency while supporting high client mobility. 
By leveraging the layered file system of Docker containers, we eliminate unnecessary transfers of a redundant and significant portion of the application file system. 
By transferring the base memory image ahead of the handoff, and transferring only the incremental memory difference when migration starts, we further reduce total transfer size. Finally, our system shows that the hand-off time is reduced by $56\%\sim 80\%$ compared to the state-of-the-art VM handoff for the edge computing platform.

% But currently, our system only supports an older version of Docker. In the future, we need to adapt the system to the newest Docker platform. 

% cite http://blog.thoward37.me/articles/where-are-docker-images-stored/ docker-0.6.7





% An example of a floating figure using the graphicx package.
% Note that \label must occur AFTER (or within) \caption.
% For figures, \caption should occur after the \includegraphics.
% Note that IEEEtran v1.7 and later has special internal code that
% is designed to preserve the operation of \label within \caption
% even when the captionsoff option is in effect. However, because
% of issues like this, it may be the safest practice to put all your
% \label just after \caption rather than within \caption{}.
%
% Reminder: the "draftcls" or "draftclsnofoot", not "draft", class
% option should be used if it is desired that the figures are to be
% displayed while in draft mode.
%
%\begin{figure}[!t]
%\centering
%\includegraphics[width=2.5in]{myfigure}
% where an .eps filename suffix will be assumed under latex, 
% and a .pdf suffix will be assumed for pdflatex; or what has been declared
% via \DeclareGraphicsExtensions.
%\caption{Simulation results for the network.}
%\label{fig_sim}
%\end{figure}

% Note that the IEEE typically puts floats only at the top, even when this
% results in a large percentage of a column being occupied by floats.


% An example of a double column floating figure using two subfigures.
% (The subfig.sty package must be loaded for this to work.)
% The subfigure \label commands are set within each subfloat command,
% and the \label for the overall figure must come after \caption.
% \hfil is used as a separator to get equal spacing.
% Watch out that the combined width of all the subfigures on a 
% line do not exceed the text width or a line break will occur.
%
%\begin{figure*}[!t]
%\centering
%\subfloat[Case I]{\includegraphics[width=2.5in]{box}%
%\label{fig_first_case}}
%\hfil
%\subfloat[Case II]{\includegraphics[width=2.5in]{box}%
%\label{fig_second_case}}
%\caption{Simulation results for the network.}
%\label{fig_sim}
%\end{figure*}
%
% Note that often IEEE papers with subfigures do not employ subfigure
% captions (using the optional argument to \subfloat[]), but instead will
% reference/describe all of them (a), (b), etc., within the main caption.
% Be aware that for subfig.sty to generate the (a), (b), etc., subfigure
% labels, the optional argument to \subfloat must be present. If a
% subcaption is not desired, just leave its contents blank,
% e.g., \subfloat[].


% An example of a floating table. Note that, for IEEE style tables, the
% \caption command should come BEFORE the table and, given that table
% captions serve much like titles, are usually capitalized except for words
% such as a, an, and, as, at, but, by, for, in, nor, of, on, or, the, to
% and up, which are usually not capitalized unless they are the first or
% last word of the caption. Table text will default to \footnotesize as
% the IEEE normally uses this smaller font for tables.
% The \label must come after \caption as always.
%
%\begin{table}[!t]
%% increase table row spacing, adjust to taste
%\renewcommand{\arraystretch}{1.3}
% if using array.sty, it might be a good idea to tweak the value of
% \extrarowheight as needed to properly center the text within the cells
%\caption{An Example of a Table}
%\label{table_example}
%\centering
%% Some packages, such as MDW tools, offer better commands for making tables
%% than the plain LaTeX2e tabular which is used here.
%\begin{tabular}{|c||c|}
%\hline
%One & Two\\
%\hline
%Three & Four\\
%\hline
%\end{tabular}
%\end{table}


% Note that the IEEE does not put floats in the very first column
% - or typically anywhere on the first page for that matter. Also,
% in-text middle ("here") positioning is typically not used, but it
% is allowed and encouraged for Computer Society conferences (but
% not Computer Society journals). Most IEEE journals/conferences use
% top floats exclusively. 
% Note that, LaTeX2e, unlike IEEE journals/conferences, places
% footnotes above bottom floats. This can be corrected via the
% \fnbelowfloat command of the stfloats package.






% conference papers do not normally have an appendix


% use section* for acknowledgment
% \section*{Acknowledgment}


% The authors would like to thank...

\bibliographystyle{IEEEtran}
% argument is your BibTeX string definitions and bibliography database(s)
\bibliography{IEEEabrv,cites}
% \bibliography{cites}
%
% <OR> manually copy in the resultant .bbl file
% set second argument of \begin to the number of references
% (used to reserve space for the reference number labels box)
% \begin{thebibliography}{1}

% \bibitem{IEEEhowto:kopka}
% H.~Kopka and P.~W. Daly, \emph{A Guide to \LaTeX}, 3rd~ed.\hskip 1em plus
%   0.5em minus 0.4em\relax Harlow, England: Addison-Wesley, 1999.

% \end{thebibliography}




% that's all folks
\end{document}


